%==============================================================================
%  MATH 231 -- Linear Algebra I (Queens College, Fall 2026)
%  EXTRA CREDIT 1 -- NumPy for Homework 2
%
%  WHAT THIS IS.  The NumPy part of Homework 2, split off as an optional
%  assignment when programming on homework became extra credit.  It carries the
%  tutorial, the two-file Python submission rules, and the nine problems that
%  were Part G, renumbered 1-9.  Homework 2 itself is now written only and never
%  mentions Python.
%
%  CROSS-REFERENCES.  Every lettered problem number here (E.1, B.2, F.6(a), ...)
%  points into homework/MATH231_homework_2.tex.  Renumber a problem there and
%  the reference here must change with it.  Where a Part G problem leaned on a
%  Homework 2 problem that was cut from the final version, it now states what it
%  needs itself:
%     Problem 3(c)  -- the equality cases of the norm inequalities (was E.9(d))
%     Problem 4(c)  -- the identity ||x^ - y^||^2 = 2 - 2cos (was E.4(d))
%     Problem 9     -- the gradient-descent step (was printed in F.11)
%  and two pieces that only checked cut problems were dropped: G.1(c) (the
%  orthogonal basis of R^4, E.7) and the D.2(c) and E.7(b) cases of G.5(b).
%
%  RUNS AHEAD.  Problem 9 uses gradient descent, veccalc S8.4, which had not been
%  lectured when this was written.  The step is printed in the problem.
%
%  NO MATRICES, with two flagged exceptions: np.corrcoef in Problem 1(b)
%  returns a small table that is read at one entry, and the optional plotting
%  section uses np.meshgrid.
%
%  THREE THINGS TO SET BEFORE HANDING THIS OUT:
%     \duedate  -- currently a placeholder.
%     \ecworth  -- how much extra credit this is worth; currently a placeholder.
%     The submission path in the "What to submit" callout, if you want a
%     different directory layout than extra-credit/ec1/ in the student repo.
%
%  Compile with pdflatex or tectonic (standard TeX Live packages only; uses
%  listings).
%==============================================================================

\documentclass[11pt]{article}

\usepackage[margin=1in]{geometry}
\usepackage{amsmath,amssymb}
\usepackage[dvipsnames]{xcolor}
\usepackage{enumitem}
\usepackage{booktabs}
\usepackage{array}
\usepackage[hidelinks]{hyperref}
\usepackage{titlesec}
\usepackage{needspace}
\usepackage{listings}

%------------------------------------------------------------------ style ----
\titleformat{\section}{\large\bfseries\sffamily\color{RoyalBlue}}{}{0em}{}
\titleformat{\subsection}{\normalsize\bfseries\sffamily}{}{0em}{}
\titlespacing*{\section}{0pt}{16pt}{7pt}

% Long \texttt names would otherwise stick out of the margin.
\setlength{\emergencystretch}{3em}

% The problem heading is a line of its own, so that part (a) -- or the opening
% sentence -- always starts below it, as in Homework 2.
\makeatletter
\newcounter{prob}
\newenvironment{prob}[1][]
  {\par\addvspace{20pt plus 6pt minus 4pt}%
   \needspace{4\baselineskip}%
   \refstepcounter{prob}%
   \noindent\textbf{Problem~\theprob}%
   \if\relax\detokenize{#1}\relax\else\ (#1)\fi\textbf{.}%
   \par\nobreak\@afterindentfalse\@afterheading}
  {\par\addvspace{20pt plus 6pt minus 4pt}}
\makeatother

% code blocks -- same style as Setup Guide 2 and Homework 1
\lstdefinestyle{term}{
  basicstyle=\ttfamily\small,
  backgroundcolor=\color{gray!10},
  frame=single, rulecolor=\color{gray!45},
  framesep=5pt, xleftmargin=8pt, xrightmargin=8pt,
  breaklines=true, columns=fullflexible, keepspaces=true,
  showstringspaces=false, aboveskip=8pt, belowskip=8pt
}
\lstdefinestyle{py}{
  style=term, language=Python,
  keywordstyle=\color{RoyalBlue}\bfseries,
  commentstyle=\color{gray!80}\itshape,
  stringstyle=\color{ForestGreen},
  emph={np,numpy}, emphstyle=\color{BrickRed}
}
\lstset{style=py}

\newcommand{\heads}[1]{\par\smallskip\noindent
  {\small\sffamily\textcolor{BrickRed}{\textbf{Watch out.}}\ #1}\par\smallskip}
\newcommand{\checkpoint}[1]{\par\smallskip\noindent
  {\small\sffamily\textcolor{ForestGreen}{\textbf{Checkpoint.}}\ #1}\par\smallskip}

% shaded call-out box.  The body is collected into a savebox first, because
% \colorbox needs its whole argument at once.
\definecolor{softbg}{HTML}{F2F6FA}
\newsavebox{\qcbox}
\newenvironment{callout}[1]
  {\par\medskip\needspace{5\baselineskip}%
   \begin{lrbox}{\qcbox}%
   \begin{minipage}{\dimexpr\textwidth-2\fboxsep-4pt}%
   \setlength{\parskip}{0.4em}%
   \textbf{\sffamily\color{RoyalBlue}#1}\par}
  {\end{minipage}\end{lrbox}%
   \begin{center}\colorbox{softbg}{\usebox{\qcbox}}\end{center}\medskip}
%--------------------------------------------------------------- notation ----
\newcommand{\R}{\mathbb{R}}
\newcommand{\vc}[1]{\mathbf{#1}}
\newcommand{\uv}[1]{\hat{\vc{#1}}}
\newcommand{\one}{\mathbf{1}}
\newcommand{\norm}[1]{\lVert #1\rVert}
\newcommand{\ip}[2]{\langle #1,#2\rangle}
\newcommand{\proj}[2]{\operatorname{proj}_{#1}(#2)}
\newcommand{\grad}{\nabla}

%--------------------------------------------------------- fill this in ------
\newcommand{\duedate}{to be announced on the class Discord}
\newcommand{\ecworth}{to be announced on the class Discord}

\title{\vspace{-1.2cm}\textbf{\sffamily MATH 231 -- Linear Algebra I}\\[4pt]
       {\large\sffamily Extra Credit 1 \textbullet\ NumPy for Homework 2}\\[2pt]
       {\normalsize\sffamily Kiely Hall 326 \textbullet\
        Instructor: Kevin Bedoya}}
\author{}
\date{}

\begin{document}
\maketitle
\vspace{-1.6cm}

\noindent\textbf{Due:} \duedate. \quad \textbf{Worth:} \ecworth.

\medskip
\noindent\textbf{This assignment is optional.} It earns extra credit and nothing
else: skipping it costs you nothing, and no later assignment depends on it. It
is the programming companion to Homework 2 --- almost every problem checks, or
takes further, a problem you worked by hand there.

\medskip
\noindent\textbf{Problem numbers.} A number with a letter --- E.1, B.2, F.6(a)
--- refers to \textbf{Homework 2}. A bare number --- Problem 3, Problem 7(c) ---
refers to this assignment.

\medskip
\noindent\textbf{Grading.} Every problem is worth the same number of points. The
assignment is graded by running \texttt{runec1.py} and reading what it prints ---
on whether it produces correct answers, not on programming style.

\medskip
\noindent\textbf{Use it to check Homework 2.} Every numerical answer on Homework 2
can be verified in a few lines of NumPy, and you are encouraged to check all of
them before submitting --- that is what the tool is for. Two limits. The hand
work on Homework 2 must still be shown: a correct number with no derivation earns
partial credit at best. And no proof can be checked this way, since code confirms
a statement only for the numbers you tried, not for all of them.

\medskip
\noindent\textbf{Collaboration} is encouraged --- work at a board together, argue
about the problems. What you submit must be written by you, in your own code.

\begin{callout}{What to submit}
Push two Python files to your own MATH 231 GitHub repository, in a folder named
\texttt{extra-credit/ec1/}:
\begin{itemize}[leftmargin=1.4em,itemsep=5pt,topsep=3pt]
  \item \texttt{ec1.py} --- your work, one function per question part;
  \item \texttt{runec1.py} --- which imports those functions and runs them all,
        labelling each answer.
\end{itemize}
Do not submit a copy of your output: \texttt{runec1.py} is \textbf{run as
submitted}, and the assignment is graded from what it prints. \emph{How to
structure your submission}, below, specifies both files in full. The rules are
the same as on Homework 1.

\textbf{Python is the only language accepted.} Any other language earns no
credit, however correct it is.

If you have not set up the repository yet, Setup Guide 1 (\S\,``Create your
folders'' and \S\,``Save your work and send it to GitHub'') is the walkthrough.
Only pushed work can be graded.
\end{callout}

\begin{callout}{Label every answer}
The labelling is built into the file structure specified under \emph{How to
structure your submission}: one function per question part in \texttt{ec1.py},
named after that part, with a comment naming it directly above; and one labelled
line per part in \texttt{runec1.py}. Code that runs correctly but cannot be
matched to a question is graded the same way an unlabelled answer sheet is.
\end{callout}

%==============================================================================
\section{What is new since Homework 1}
%==============================================================================
\noindent This assignment assumes Homework 1's Part F: arrays, \texttt{+},
\texttt{-}, scaling, \texttt{np.linalg.norm}, \texttt{np.isclose},
\texttt{np.allclose}, and the two-file submission. The short tutorial below
covers only what is new --- vectors of any length, the dot product, the other
norms, angles, and functions of one and two variables. Then comes the submission
specification (the same rules as before, with the file names changed), then
Problems 1--9. The \textbf{Python Recipe Book} (\texttt{lectures/python-recipe/})
covers the vector calls in its \S2 and floating-point behaviour in its \S7; the
remaining calls are introduced below or in the problem that uses them.

\subsection{Vectors in \texorpdfstring{$\R^{n}$}{Rn}}

Nothing changes when a vector has more than two entries.

\begin{lstlisting}
import numpy as np

u = np.array([1.0, -2.0, 0.0, 3.0])   # a vector in R^4
len(u)                                # 4 -- the n in R^n

np.zeros(5)                           # array([0., 0., 0., 0., 0.])
np.ones(5)                            # the all-ones vector 1 in R^5
np.arange(1, 6)                       # array([1, 2, 3, 4, 5])

e3 = np.zeros(5)                      # the standard basis vector e_3 in R^5:
e3[2] = 1.0                           # slot 3 is index 2
\end{lstlisting}

\heads{Adding vectors of different lengths raises an error, exactly as in
lecture: \texttt{np.ones(3) + np.ones(4)} fails with \texttt{ValueError}.}

\subsection{The dot product}

\begin{lstlisting}
v = np.array([2.0, 1.0, -1.0, 0.0])

u @ v              # 0.0  -- the dot product <u, v>
np.dot(u, v)       # 0.0  -- the same number
u @ u              # 14.0 -- <u, u> = ||u||^2
\end{lstlisting}

\heads{\texttt{u * v} is still \emph{not} the dot product. It multiplies entry by
entry and returns a vector. \texttt{np.sum(u * v)} does equal \texttt{u @ v},
but write \texttt{u @ v}: it says what you mean.}

\subsection{The other norms, and a few reductions}

\begin{lstlisting}
np.linalg.norm(u)               # 3.7416...  ||u||_2 = sqrt(14)
np.linalg.norm(u, ord=1)        # 6.0        ||u||_1
np.linalg.norm(u, ord=np.inf)   # 3.0        ||u||_inf
np.linalg.norm(u, ord=3)        # ||u||_3

np.abs(u)                       # array([1., 2., 0., 3.]) -- entrywise |u_i|
np.sum(u)                       # 2.0  -- the sum of the entries
np.max(np.abs(u))               # 3.0  -- the largest |u_i|
np.mean(u)                      # 0.5  -- the average entry
\end{lstlisting}

\subsection{Angles}

\begin{lstlisting}
c = (u @ v) / (np.linalg.norm(u) * np.linalg.norm(v))  # cos(u, v)
theta = np.arccos(np.clip(c, -1.0, 1.0))                # radians, in [0, pi]
np.degrees(theta)                                       # 90.0
\end{lstlisting}

\noindent \texttt{np.clip} guards against rounding pushing \texttt{c} a hair
outside $[-1,1]$, where \texttt{np.arccos} is undefined and returns
\texttt{nan}.

\subsection{Functions of one and two variables}

A mathematical function becomes a Python function with \texttt{def}. NumPy
supplies the elementary functions.

\begin{lstlisting}
def f(x, y):
    return x**3 - 3*x*y + y**2

f(1.0, 2.0)                  # -1.0

np.exp(1.0)                  # 2.718281828...   e^1
np.log(2.0)                  # 0.693147...      the NATURAL log, ln 2
np.sin(np.pi / 6)            # 0.4999999...     sin(pi/6)
np.sqrt(9.2)                 # 3.03315...
\end{lstlisting}

\heads{\texttt{np.log} is $\ln$. The base-$10$ logarithm is \texttt{np.log10};
you will not need it.}

\noindent A function can be handed to another function, just like a number. That
is how Problem 7 differentiates \emph{any} $f$ with one piece of code:

\begin{lstlisting}
def twice(g, x):
    return 2 * g(x)

twice(np.exp, 0.0)           # 2.0
\end{lstlisting}

\subsection{Loops over step sizes, and random vectors}

\begin{lstlisting}
hs = 10.0 ** -np.arange(1, 16)     # 1e-1, 1e-2, ..., 1e-15
for h in hs:
    print(h, (np.exp(h) - 1) / h)

rng = np.random.default_rng(231)   # a fixed seed, as in Homework 1
x = rng.normal(size=10)            # a random vector in R^10
\end{lstlisting}

%==============================================================================
\section{How to structure your submission}
%==============================================================================
The assignment is submitted as \textbf{two} Python files with different jobs,
exactly as on Homework 1.

\begin{itemize}[leftmargin=1.4em,itemsep=6pt,topsep=4pt]
  \item \texttt{ec1.py} --- \textbf{your work}. One function per lettered part,
        named \texttt{q}\,\emph{number}\,\emph{letter}: Problem 2(a) becomes
        \texttt{q2a}, Problem 7(c) becomes \texttt{q7c}. Each function
        \textbf{returns} its answer; when an answer has several pieces, return
        them together as a tuple or a list. Helper functions (\texttt{my\_dot},
        \texttt{pnorm}, \texttt{proj}, \dots) are defined in this file too, and
        may be shared between parts.
  \item \texttt{runec1.py} --- \textbf{the driver}. It imports those functions
        from \texttt{ec1.py} and calls them in order, one line per part,
        printing a label before each answer. It contains no mathematics of its
        own. \textbf{This is the file that is run to grade the assignment.}
\end{itemize}

\medskip
\noindent\textbf{\texttt{ec1.py} --- definitions only.} Nothing may run when the
file is imported: define your functions and stop. A \texttt{print} sitting at the
left margin fires the moment \texttt{runec1.py} imports the file, and its output
lands above the labels, where it cannot be matched to a question.

\begin{lstlisting}
# ec1.py -- the work.  Definitions only; nothing runs on import.

import numpy as np

def pnorm(v, p):                    # a helper, used by several parts
    ...

# --- 2(c) ---------------------------------------------------------
def q2c():
    v = np.array([-3.0, 4.0])
    return [pnorm(v, p) for p in (1, 2, 3, 4, 10, 100)] + [infnorm(v)]
\end{lstlisting}

\noindent\textbf{\texttt{runec1.py} --- one line per part.} Import the functions,
then call them in the order the questions appear. Each line prints a label
followed by the answer that call produced:

\begin{lstlisting}
# runec1.py -- runs every question in order.  No mathematics in this file.

import ec1

print("Output of question 1 part a : ", ec1.q1a())
print("Output of question 1 part b : ", ec1.q1b())
print("Output of question 2 part a : ", ec1.q2a())
\end{lstlisting}

\medskip
\noindent\textbf{The label is required, and its wording is fixed:}

\begin{center}
\texttt{"Output of question <number> part <letter> : "}
\end{center}

\noindent So Problem 7(c) is labelled \texttt{Output of question 7 part c :}\ .

\heads{Several parts below ask you to ``answer in a comment''. Those comments go
in \texttt{ec1.py}, directly under the function for that part, and are read by
the grader. A part that asks \emph{only} for a comment still gets its function
and its labelled line: define it to return the string \texttt{"see comment"}.
Everything else is \emph{returned}, not printed: \texttt{runec1.py} does all the
printing.}

\heads{Problems 2, 3, 4, 6 and 8 use random numbers. Create the generator
\emph{inside} the function that uses it, with the seed given, so that every run
of \texttt{runec1.py} prints the same thing.}

\checkpoint{Run \texttt{runec1.py} once before you push and read the output top
to bottom. Every part of every question should appear, in order, under its own
label, with no error and nothing unlabelled.}

%==============================================================================
\section{The problems}
%==============================================================================

\begin{prob}[Checking Homework 2, Part E]
\begin{enumerate}[label=\textbf{(\alph*)},leftmargin=2.6em,itemsep=7pt,topsep=6pt]
  \item Enter $\vc{u}$, $\vc{v}$, $\vc{w}$ of Problem E.1 and return the values
        asked for in every one of its parts (a)--(g), in order, as one list.
        The angle in (e) should be in degrees.
  \item For the data $\vc{x}$ and $\vc{y}$ of Problem E.5, return:
        $\proj{\one}{\vc{x}}$; the deviation vector; its dot product with
        $\one$; $\sigma^{2}$ and $\sigma$; and $r$. Then return, alongside your
        own values, \texttt{np.mean(x)}, \texttt{np.var(x)}, \texttt{np.std(x)}
        and \texttt{np.corrcoef(x, y)[0, 1]}, and check each pair with
        \texttt{np.isclose}. (\texttt{np.corrcoef} returns a small table of
        correlations rather than a single number; its entry \texttt{[0, 1]} is
        $r$. That indexing is all you need to know about it.)
\end{enumerate}

\noindent State in a comment under each function whether every value agrees with
your hand computation. If any disagree, find the error --- it is far more often
in the hand computation than in the code.
\end{prob}

\begin{prob}[Writing the dot product and the $p$-norm yourself]
\begin{enumerate}[label=\textbf{(\alph*)},leftmargin=2.6em,itemsep=7pt,topsep=6pt]
  \item Write \texttt{my\_dot(u, v)} using a \texttt{for} loop over the indices
        --- not \texttt{@}, \texttt{np.dot}, or \texttt{np.sum(u * v)}. It
        should raise a \texttt{ValueError} if the lengths differ. Return
        \texttt{my\_dot} on the three pairs of Problem E.1(d), each checked
        against \texttt{@} with \texttt{np.isclose}.
  \item Write \texttt{pnorm(v, p)} for $p\ge1$ and \texttt{infnorm(v)}, using
        only \texttt{np.abs}, \texttt{np.sum}, \texttt{np.max} and
        \texttt{**}. With \texttt{rng = np.random.default\_rng(231)}, draw five
        random vectors in $\R^{7}$ and, for each, compare \texttt{pnorm} with
        \texttt{np.linalg.norm(v, ord=p)} for $p=1,2,3,4.5$, and
        \texttt{infnorm} with \texttt{ord=np.inf}. Return \texttt{True} only if
        all $25$ comparisons pass.
  \item For $\vc{v}=[\,-3,\,4\,]$, return \texttt{pnorm(v, p)} for
        $p=1,2,3,4,10,100$, and \texttt{infnorm(v)}. Compare with Problem
        B.1(iii).
  \item With \texttt{rng = np.random.default\_rng(7)}, draw one random vector in
        $\R^{10}$ and return \texttt{pnorm} for $p=1,2,4,8,16,32,64,128$,
        followed by \texttt{infnorm}. In a comment: at roughly which $p$ is
        $\norm{\vc{v}}_p$ within $1\%$ of $\norm{\vc{v}}_\infty$?
\end{enumerate}
\end{prob}

\begin{prob}[The norm inequalities at scale]
Problem E.7 proved, for $\vc{v}\in\R^{n}$,
\[
  \norm{\vc{v}}_\infty\le\norm{\vc{v}}_2\le\norm{\vc{v}}_1
  \le\sqrt{n}\,\norm{\vc{v}}_2\le n\,\norm{\vc{v}}_\infty .
\]
\begin{enumerate}[label=\textbf{(\alph*)},leftmargin=2.6em,itemsep=7pt,topsep=6pt]
  \item With \texttt{rng = np.random.default\_rng(231)}, draw $10{,}000$ random
        vectors in $\R^{10}$ and count how many satisfy all four inequalities.
  \item In the same run, return the largest and the smallest value of
        $\norm{\vc{v}}_1/\norm{\vc{v}}_2$ you observed, next to the bounds $1$
        and $\sqrt{10}$ that E.7 predicts.
  \item Find a nonzero vector in $\R^{10}$ for which
        $\norm{\vc{v}}_2\le\norm{\vc{v}}_1$ is an equality, and one for which
        $\norm{\vc{v}}_1\le\sqrt{n}\,\norm{\vc{v}}_2$ is an equality. Build both
        in code, and return \texttt{np.isclose} verdicts confirming it.
  \item Answer in a comment. The random vectors came much closer to one of the
        two bounds in (b) than to the other. Which one? Using your vectors from
        (c), say what a vector must look like to attain each bound, and why one
        kind is far rarer among random draws than the other. Finally: does the
        count in (a) prove E.7?
\end{enumerate}
\end{prob}

\begin{prob}[Similarity and nearest neighbours]
\begin{enumerate}[label=\textbf{(\alph*)},leftmargin=2.6em,itemsep=7pt,topsep=6pt]
  \item Write \texttt{cos\_sim(u, v)}. For the documents of Problem E.4, return
        $\cos(A,X)$ and $d_2(A,X)$ for $X=B,C,D$.
  \item For the depot problem B.2, return, for each of \texttt{ord=1},
        \texttt{ord=2} and \texttt{ord=np.inf}, the name (\texttt{"A"},
        \texttt{"B"} or \texttt{"C"}) of the nearest customer. Let the code find
        the minimum; do not type the answer in.
  \item For nonzero $\vc{x},\vc{y}\in\R^{n}$, the unit vectors
        $\uv{x}=\vc{x}/\norm{\vc{x}}$ and $\uv{y}=\vc{y}/\norm{\vc{y}}$ satisfy
        \[
          \norm{\uv{x}-\uv{y}}^{2} \;=\; 2-2\cos(\vc{x},\vc{y}) ,
        \]
        so once every document is normalized, ranking by distance and ranking
        by cosine similarity give the same order. With
        \texttt{rng = np.random.default\_rng(231)}, draw $1000$ pairs of random
        vectors in $\R^{50}$ and check the identity on every pair with
        \texttt{np.isclose}. Return the number of pairs that pass.
\end{enumerate}
\end{prob}

\begin{prob}[Projections and coordinates in $\R^{n}$]
\begin{enumerate}[label=\textbf{(\alph*)},leftmargin=2.6em,itemsep=7pt,topsep=6pt]
  \item Write \texttt{proj(u, v)}, returning $\proj{\vc{u}}{\vc{v}}$. Use it on
        Problem A.3(a) and on the pair of Problem D.3, and return each
        projection together with the dot product of its orthogonal piece against
        the vector projected onto.
  \item Write \texttt{coords(basis, x)}, where \texttt{basis} is a Python list of
        mutually orthogonal nonzero arrays, returning the list of coefficients
        $\ip{\vc{b}}{\vc{x}}/\ip{\vc{b}}{\vc{b}}$. Write also
        \texttt{is\_orthogonal\_set(vectors)}, which checks every pair with
        \texttt{np.isclose}. Return the coordinates (and the orthogonality
        verdict) for every vector in Problems D.1(a) and D.1(c).
  \item For each case in (b), rebuild the vector from its coordinates and return
        the \texttt{np.allclose} verdict against the original.
  \item Now misuse \texttt{coords} on the basis $\{[2,1],[1,3]\}$ of Problem
        C.1, which is not orthogonal, with $\vc{x}=[\,5,\,5\,]$. Return the
        coefficients it produces, the vector they rebuild, and the
        \texttt{np.allclose} verdict. Explain in a comment which step of the
        derivation in vector unit \S21.6 fails when the basis is not
        orthogonal.
\end{enumerate}
\end{prob}

\begin{prob}[Life in high dimensions]
\begin{enumerate}[label=\textbf{(\alph*)},leftmargin=2.6em,itemsep=7pt,topsep=6pt]
  \item Return the angle in degrees between $\one$ and $\vc{e}_1$ in $\R^{n}$
        for $n=2,3,4,10,100,10000$, computed from the vectors themselves (not
        from the formula $1/\sqrt{n}$ of Problem E.6). Compare with E.6(b).
  \item With \texttt{rng = np.random.default\_rng(231)}, for each
        $n=2,10,100,1000,10000$ draw $2000$ pairs of random vectors in $\R^{n}$
        and compute the average of $\lvert\cos(\vc{x},\vc{y})\rvert$ over the
        pairs. Return the five averages, and the five values of
        $\sqrt{n}$ times the average.
  \item Answer in a comment, in two or three sentences: what happens to the angle
        between two random vectors as $n$ grows? Search engines and language
        models represent words and documents as vectors in $\R^{n}$ with $n$ in
        the hundreds or thousands and entries of both signs. If two such vectors
        have cosine similarity $0.3$, is that a large number or a small one, and
        why? (Raw word counts, as in Problem E.4, behave differently: their
        entries are never negative, so their cosine is never negative, and the
        experiment in (b) does not describe them.)
\end{enumerate}
\end{prob}

\begin{prob}[The derivative as a limit, numerically]
\begin{enumerate}[label=\textbf{(\alph*)},leftmargin=2.6em,itemsep=7pt,topsep=6pt]
  \item Write \texttt{diff\_quotient(f, x, h)}, returning
        $\bigl(f(x+h)-f(x)\bigr)/h$. For $f=\exp$ at $x=1$, return the error
        $\lvert\text{quotient}-e\rvert$ for each $h=10^{-1},10^{-2},\dots,10^{-15}$,
        and the value of $h$ that gave the smallest error.
  \item Answer in a comment: the error shrinks as $h$ shrinks, and then grows
        again. Explain both halves. (The Recipe Book's ``Watching
        cancellation'', \S7, is the second half.)
  \item Return $(e^{h}-1)/h$ and $(1+h)^{1/h}$ for $h=10^{-1},\dots,10^{-8}$.
        Which two numbers from the vector calculus lecture are they approaching?
  \item With $h=10^{-6}$, use \texttt{diff\_quotient} to check your answers to
        Problem F.1(a), (b), (e) and (f) at $x=2$: code your hand formula for
        each derivative, and compare with \texttt{np.isclose(..., rtol=1e-4)}.
        Return the four verdicts.
\end{enumerate}
\end{prob}

\begin{prob}[Partial derivatives and the gradient, numerically]
\begin{enumerate}[label=\textbf{(\alph*)},leftmargin=2.6em,itemsep=7pt,topsep=6pt]
  \item Write \texttt{partial\_x(f, x, y, h=1e-6)} and
        \texttt{partial\_y(f, x, y, h=1e-6)} from the limit definitions, and
        \texttt{grad(f, x, y)} returning the gradient as a NumPy array. Return
        \texttt{grad} of $x^{3}-3xy+y^{2}$ at $(1,2)$, $(2,-1)$ and $(0,0)$, next
        to your hand values from Problem F.6(a), with \texttt{np.allclose}
        verdicts (use \texttt{atol=1e-4}).
  \item Return $\norm{\grad f(2,-1)}$ computed numerically, next to the exact
        value.
  \item Check your answers to Problem F.4(b), (c) and (h) at the point
        $(1,\,0.5)$: code your hand formulas and compare. Return the verdicts.
  \item For $f(x,y)=\sqrt{x^{2}+y^{2}}$, with
        \texttt{rng = np.random.default\_rng(231)}, draw five random points and
        return $\norm{\grad f}$ at each. Compare with Problem F.6(d).
\end{enumerate}
\end{prob}

\begin{prob}[Running ahead: gradient descent]
This problem uses \S8.4 of the vector calculus unit, which has not been lectured
yet. Everything you need is here:

\begin{callout}{From \S8.4}
\textbf{Gradient descent} minimizes a function $f:\R^{2}\to\R$ by repeatedly
stepping in the direction of steepest descent, $-\grad f$:
\[
  \vc{x}_{k+1} \;=\; \vc{x}_k-\alpha\,\grad f(\vc{x}_k),
\]
with a fixed step size $\alpha>0$ and a chosen starting point $\vc{x}_0$.
\end{callout}

\begin{enumerate}[label=\textbf{(\alph*)},leftmargin=2.6em,itemsep=7pt,topsep=6pt]
  \item Write \texttt{descend(grad\_f, x0, alpha, steps)}, returning the list of
        all points $\vc{x}_0,\vc{x}_1,\dots,\vc{x}_{\text{steps}}$. For
        $f(x,y)=x^{2}+4y^{2}$, with $\vc{x}_0=(2,1)$ and $\alpha=0.1$, return
        $\vc{x}_1$ and $\vc{x}_2$, and the point and the value of $f$ after $50$
        steps.
  \item Repeat with $\alpha=0.25$ and with $\alpha=0.3$. Return the point and
        the value of $f$ after $50$ steps for each.
  \item Answer in a comment. Show by hand that for this $f$ one step is
        $x_{k+1}=(1-2\alpha)x_k$ and $y_{k+1}=(1-8\alpha)y_k$. Use it to explain
        all three runs, and find the range of $\alpha$ for which the iteration
        converges to the minimum.
\end{enumerate}

\noindent Finally, answer in a comment at the end of your file, in three or four
sentences: which problems on Homework 2 could NumPy check for you, and which
could it not? Name at least one proof problem and say precisely what a computer
run would and would not establish about it.
\end{prob}

\begin{callout}{What NumPy can and cannot do for you}
\renewcommand{\arraystretch}{1.25}
\begin{tabular}{@{}l >{\raggedright\arraybackslash}p{10.4cm}@{}}
\toprule
\textbf{Homework 2} & \textbf{Can NumPy check it?} \\
\midrule
A --- projection         & A.1--A.3 and the numbers in A.4--A.6, yes; the proofs, no. \\
B --- norms and metrics  & B.1 and B.2, yes; B.3--B.5, no. \\
C --- bases and span     & C.1--C.3 and C.7, yes; the proofs in C.4--C.6, no. \\
D --- orthogonality      & D.1, D.2 and the numbers in D.3, yes; the proof in D.3, D.4, D.5, no. \\
E --- $\R^{n}$           & The numbers, yes; E.2, E.3(c), E.5(a), E.7, no. \\
F --- derivatives        & Every derivative at a point, yes; F.3, F.6(c)--(d), no. \\
\midrule
This assignment          & It \emph{is} the check. \\
\bottomrule
\end{tabular}

\smallskip
A numerical derivative checks a formula at the points you try. It cannot tell you
the formula is right everywhere --- but a mismatch at one point proves the
formula wrong.
\end{callout}

%==============================================================================
\section{Optional and ungraded: drawing the pictures}
%==============================================================================
\noindent Not required, not graded, and not needed for any later assignment. It
needs \texttt{matplotlib}, as in Homework 1:

\begin{lstlisting}
python -m pip install matplotlib
\end{lstlisting}

\heads{Keep this out of \texttt{ec1.py} and \texttt{runec1.py}, in a scratch file
of its own. \texttt{plt.show()} halts a program until you close the window, so a
plot left in your submission stalls the run it is graded from.}

\noindent The unit balls of vector unit \S14.8, for several values of $p$. Each
point of the ordinary unit circle is rescaled to have $p$-norm $1$:

\begin{lstlisting}
import numpy as np
import matplotlib.pyplot as plt

t = np.linspace(0, 2*np.pi, 400)
x, y = np.cos(t), np.sin(t)
for p in [1, 1.5, 2, 4, 10]:
    size = (np.abs(x)**p + np.abs(y)**p) ** (1/p)   # the p-norm of each point
    plt.plot(x / size, y / size, label=f"p = {p}")

plt.gca().set_aspect("equal")
plt.legend()
plt.grid(alpha=0.3)
plt.show()
\end{lstlisting}

\noindent The level curves of $g(x,y)=x^{2}+4y^{2}$ --- the function Problem 9
minimizes --- with its gradient drawn at a grid of points. \texttt{np.meshgrid}
builds the grid of points; you do not need to understand the shape of what it
returns to use it.

\begin{lstlisting}
xs = np.linspace(-2, 2, 200)
X, Y = np.meshgrid(xs, xs)
plt.contour(X, Y, X**2 + 4*Y**2, levels=12)

xg = np.linspace(-2, 2, 9)
XG, YG = np.meshgrid(xg, xg)
plt.quiver(XG, YG, 2*XG, 8*YG, color="C3")   # the gradient [2x, 8y]

plt.gca().set_aspect("equal")
plt.show()
\end{lstlisting}

\noindent Notice that the arrows cross every level curve at a right angle, and
that they do not point away from the centre except along the two axes.

\bigskip
\noindent\rule{\textwidth}{0.4pt}

\smallskip
\noindent{\small Questions about any problem here belong in the class Discord,
where everyone can see the answer. Questions that are specific to your own work
--- ``why does my \texttt{runec1.py} stop at Problem 4?'' --- are better brought
to office hours, Thursdays 3:00--4:00 PM in Kiely Hall 508. Please email ahead.}

\end{document}
